% Synthetic skill-evaluation manuscript; not research findings.
% Derived from the supplied multi-file ICML project. Official files unchanged.
\documentclass[journal]{IEEEtai}
\usepackage{microtype,graphicx,booktabs,hyperref}
\usepackage{algorithm,algorithmic}
\usepackage{amsmath,amssymb}
\newtheorem{PWTheorem}{Theorem}
\newcommand{\riskname}{selective risk}
\def\methodname{Gate}
\providecommand{\theHalgorithm}{\arabic{algorithm}}
\hypersetup{hidelinks,pdftitle={A Fixed Confidence Gate for Selective Classification},
  pdfauthor={},pdfsubject={Synthetic skill evaluation},pdfkeywords={selective classification, confidence threshold, coverage, selective risk}}
\title{A Fixed Confidence Gate for Selective Classification}
\author{Anonymous submission}
\markboth{Fixed Confidence Gate}{Fixed Confidence Gate}
\begin{document}
\maketitle
\begin{abstract}
Selective classification allows a model to abstain from uncertain predictions,
but accuracy among accepted predictions describes a selected population.
This manuscript examines a fixed confidence gate in a synthetic evaluation
fixture, with coverage reported alongside conditional accuracy. The gate
accepts predictions at or above a specified confidence threshold and otherwise
abstains; the comparator uses the same classifier and accepts every prediction.
Neither method retrains the model. Across two synthetic datasets, conditional
accuracy increases on one and decreases on the other. These observations do
not establish an overall accuracy improvement because the evaluation
populations differ. A range bound places selective risk between its stated
limits under bounded loss and positive coverage; it establishes neither a
generalization guarantee nor superiority over the comparator. The fixture
illustrates the need to interpret abstention metrics with their coverage and
assumptions. Its values are not research findings. Statistical uncertainty,
runtime, distribution shift, and deployment outcomes remain unevaluated,
and missing classifier and data-generation details limit reproducibility.
\end{abstract}
\begin{IEEEImpStatement}
Abstention changes which predictions contribute to a reported accuracy.
Researchers and practitioners therefore need coverage to interpret accuracy
among accepted predictions. This synthetic fixture illustrates that reporting
only conditional accuracy can obscure the population being evaluated. It
provides an example for checking whether manuscript language preserves this
distinction during editing and template migration. The fixture is not evidence
of a deployable classification system or an improvement in decision quality.
No runtime, distribution-shift, or deployment assessment has been performed,
and statistical uncertainty is unavailable. Any proposed practical use would
require reproducible model and data descriptions, suitable comparisons,
and evaluation of the relevant operating conditions. The numerical values
remain synthetic inputs to a skill evaluation, rather than research findings
or demonstrated social benefits.
\end{IEEEImpStatement}
\begin{IEEEkeywords}
Confidence threshold, coverage, selective classification, selective risk
\end{IEEEkeywords}
\section{Introduction}
Selective classification allows a model to abstain from uncertain predictions.
We evaluate a fixed confidence gate using synthetic fixtures. Because the gate
selects which predictions are evaluated, we report coverage alongside accuracy
among accepted predictions.
\section{Method}\label{sec:method}
Gate uses a fixed classifier and accepts a prediction when its maximum class
probability is at least 0.8; otherwise, it abstains. Always uses the same classifier
and accepts every prediction. Neither method retrains the model.
For a fixed classifier $h$ and acceptance function $a(x)\in\{0,1\}$, assume
$0\leq\ell(h(x),y)\leq1$ and positive coverage $C=\mathbb{E}[a(X)]>0$.
Selective risk is
\begin{equation}\label{eq:risk}
R_{\mathrm{sel}}=\frac{\mathbb{E}[a(X)\ell(h(X),Y)]}{C}.
\end{equation}
\begin{PWTheorem}\label{thm:bounded}
Under these assumptions, $0\leq R_{\mathrm{sel}}\leq1$.
\end{PWTheorem}
This range bound provides neither a generalization guarantee nor a comparison
between Gate and Always. Appendix~\ref{app:proof} gives the full proof.
\section{Synthetic Evaluation}\label{sec:results}
Alpha and Beta each contain 1000 synthetic samples. Both methods use the same
model and 5 evaluation seeds; results are seed means, with no variance estimates.

\begin{table}[htbp]\centering
\caption{Synthetic fixture values. Always accuracy uses all samples;
Gate accuracy is conditional on acceptance. These values are not research findings.}
\label{tab:results}
\begin{tabular}{lrrr}\hline
Dataset & Always (\%) & Gate (\%) & Coverage (\%)\\\hline
Alpha & 80.0 & 84.0 & 60.0\\
Beta & 76.0 & 75.5 & 58.0\\\hline
\end{tabular}\end{table}

Table~\ref{tab:results} shows a higher conditional accuracy on Alpha and a lower
conditional accuracy on Beta. Different evaluation populations prevent interpreting
these differences as an overall accuracy improvement. No significance tests,
runtime measurements, or distribution-shift evaluations were performed.
\section{Conclusion and Future Work}
The synthetic evaluation compares a fixed gate with a single baseline. It does
not establish state-of-the-art performance or a guarantee of lower risk.
The classifier and data-generation details remain missing and are required for
reproducibility. Conditional accuracy and coverage describe the accepted
population; deployment benefits remain untested.
Future work should assess uncertainty, runtime, and distribution shift after
the reproducibility gaps are resolved. The source draft used the ICML author guidelines~\cite{formatguide};
this version uses the locally verified IEEE TAI template.
\label{pw:main-end}

\bibliographystyle{IEEEtran}\bibliography{references}
\appendices
\section{Proof of the Range Bound}\label{app:proof}
Since $0\leq a(X)\ell(h(X),Y)\leq a(X)$, taking expectations gives
$0\leq\mathbb{E}[a(X)\ell(h(X),Y)]\leq C$. Division by $C>0$ proves
Theorem~\ref{thm:bounded}. This argument applies to the specified distribution;
it does not establish a bound under distribution shift.


\end{document}
